% Traducción revisada sobre tex/en/ch01.tex y el borrador de Anahi.
% Los números de párrafo remiten al DOCX inglés; véase notes/ch01-summary.md.

% DOCX paragraph 55
\chapter{La gravedad… ¡atrae!}\label{ch:01}

% DOCX paragraph 58
Como ya se ha dicho, el título de este capítulo pretende ser menos trivial de lo que
pudiera parecer.

% DOCX paragraph 59
En esta ocasión no quiero limitarme a subrayar la existencia de una «fuerza» física
ejercida entre dos objetos con masa, la cual hace que estos se atraigan entre sí incluso
cuando se encuentran a gran distancia. De hecho, en el capítulo 3 descubriremos que
esta idea, por más extendida que esté y por fácil de entender que sea, es incorrecta y,
al menos en parte, engañosa. Con este título quiero hacer hincapié en que hay algo
---la gravedad--- que nos atrae en un sentido amplio y, por encima de todo, atrae
nuestra atención. Además de actuar sobre los objetos físicos, atenaza nuestra
imaginación de manera irresistible; es capaz de impulsar nuestros pensamientos hacia
horizontes radicalmente nuevos, que en nada se parecen a los habituales, y exponernos
a escenarios que se extienden mucho más allá de nuestra experiencia cotidiana.

% DOCX paragraph 60
Pero vayamos paso a paso. Para definir mejor qué es la gravedad, resulta útil dividir
nuestro conocimiento en tres niveles distintos, pero conectados entre sí. En particular,
podríamos decir que tenemos un \emph{conocimiento instintivo}\index{conocimiento!instintivo},
un \emph{conocimiento racional}\index{conocimiento!racional} y, por último, un
\emph{conocimiento imaginativo}\index{conocimiento!imaginativo} de la gravedad.

% DOCX paragraph 61
Veamos en qué consisten y cómo difieren entre sí.

% DOCX paragraph 64
\section{Conocimiento instintivo}\label{sec:01-instinctive}

% DOCX paragraph 66
Todos sabemos que el instinto es la tendencia natural e intrínseca ---es decir, que no
requiere la intervención del razonamiento o la reflexión--- a adoptar un comportamiento
particular. Un ejemplo de instinto sería aquello que hace que nos sobresaltemos y nos
protejamos la cabeza cuando nos sorprende un sonido repentino, fuerte y desconocido.
Expresado en estos términos, es difícil creer que algo instintivo o irracional pueda
ligarnos a la gravedad. Y, sin embargo, es así.

% DOCX paragraph 67
Cualquier persona que haya tenido en brazos a un recién nacido sin duda habrá sido
testigo de primera mano del llamado reflejo de Moro\index{Moro!reflejo}, que recibe su
nombre del pediatra austríaco Ernst Moro\index{Moro!Ernst} (1874-1951). Es uno de los
principales reflejos neonatales, muy utilizado para evaluar el funcionamiento del
sistema nervioso central. Para inducir este reflejo, basta con tomar a un recién nacido,
incluso a los pocos segundos de su nacimiento, elevarlo en posición horizontal y
crearle la sensación de caída libre\index{caida libre@caída libre}, asegurándose, por
supuesto, de que no haya ningún peligro. La respuesta del bebé a esta inesperada y
aparente falta de sujeción es una reacción de sorpresa, acompañada de una repentina
apertura de brazos y manos en busca de un asidero, como se muestra en la
figura~\ref{fig:01-01}.

% DOCX paragraph 68
Desde un punto de vista médico, la aparición de este reflejo es una prueba importante
del perfecto funcionamiento fisiológico del sistema nervioso central de la criatura.
Por eso a todos nos han sometido a esa prueba. Cualquier padre o madre que la haya
presenciado ---en mi caso han sido tres veces--- sabe muy bien el alivio que sintió
al ver que su hijo o hija respondía de manera correcta a este extraño estímulo.

% DOCX paragraph 70
\begin{figure}[htbp]
\centering
\reflectbox{\includegraphics[width=0.65\linewidth]{../figs/source/en/fig-01-01.jpg}}

% DOCX paragraph 71
\caption{\emph{Ejemplo del reflejo de Moro: una recién nacida reacciona de forma
instintiva a la falta de sujeción abriendo los brazos e intentando agarrarse a algo
para detener la caída. © V. Tverdokhlib/depositphotos.com.}}\label{fig:01-01}
\end{figure}

% DOCX paragraph 73
Desde una perspectiva antropológica, el reflejo de Moro nos recuerda nuestro pasado
como primates, cuando, con toda probabilidad, teníamos que estar listos enseguida para
seguir a nuestras madres, bien agarrados a su espalda. De todos modos, lo que más nos
interesa aquí es su significado físico. La presencia de este
instinto\index{conocimiento!instintivo} unos pocos segundos después del nacimiento
---cuando estamos completamente indefensos y no sabemos nada del mundo que nos rodea---
nos revela una importante verdad sobre nuestra interacción con la gravedad: la
conocemos instintivamente, mucho antes de tener interacciones conscientes con el resto
del universo físico. O sea que, después de haber pasado nueve cómodos meses en el seno
de nuestras madres, esencialmente aislados de todo, somos capaces de responder
inmediatamente a la gravedad o, más bien, a su ausencia. No es un detalle menor.

% DOCX paragraph 74
El reflejo de Moro desaparece aproximadamente a los seis meses, de manera que nuestro
conocimiento de la gravedad, si bien en parte sigue siendo instintivo, se va modificando
con el tiempo a medida que desarrollamos la capacidad de observar el universo físico
y comprender sus leyes.

% DOCX paragraph 77
\section{Conocimiento racional}\label{sec:01-rational}

% DOCX paragraph 79
A medida que nuestra experiencia del mundo se amplía y se pulen nuestras habilidades
intelectuales, nuestro conocimiento de la gravedad pasa de un nivel instintivo a uno
racional. En otras palabras, se convierte en una parte integral de nuestras expectativas
sobre cómo funciona el mundo a nuestro alrededor. Esto queda claramente demostrado
mediante sencillos experimentos visuales que se llevan a cabo con niños de corta edad
empleando dibujos animados. Aunque algunos de los sujetos que participan todavía no
son capaces de caminar, todos pueden interpretar el movimiento de un objeto coherente
con la presencia de un campo gravitatorio. El ejemplo clásico es el de una esfera
rodando por una mesa: los niños reaccionan de manera distinta, en sus expresiones
faciales y en los movimientos de los ojos, dependiendo de que la esfera, al llegar al
borde de la mesa, caiga, continúe su movimiento aunque no tenga sujeción o bien, lo
cual es más raro aún, empiece a volar. El hecho de que estos niños tan pequeños sepan
cuál debería ser el comportamiento correcto de la pelota vuelve a confirmar hasta qué
punto el conocimiento de la «fuerza» de la gravedad está arraigado en nuestra mente.

% DOCX paragraph 80
Este conocimiento desempeña una función fundamental en nuestra percepción\index{percepcion@percepción}
racional de la realidad. Gracias a esta profunda influencia, nuestros cerebros son
capaces de resolver problemas dinámicos de gran complejidad en un tiempo muy breve y
prácticamente sin esfuerzo.\index{conocimiento!racional} Un ejemplo sencillo es la
acción de bajar unas escaleras a toda prisa: es uno de los problemas más complicados
a la hora de programar robots, y algo en lo que a menudo fracasan; no obstante, los
humanos nos desenvolvemos muy bien sin siquiera pensarlo conscientemente. Y, sin
embargo, establecer en qué secuencia y a qué velocidad debemos llevar a cabo nuestros
movimientos para asegurar el sutil equilibrio entre las distintas fuerzas que compiten
con la fuerza gravitatoria no es trivial en absoluto.

% DOCX paragraph 81
Por último, hay otra propiedad de la gravedad que merece nuestra consideración: su
capacidad de estimular nuestra imaginación.

% DOCX paragraph 84
\section{Conocimiento imaginativo}\label{sec:01-imaginative}

% DOCX paragraph 86
Si está claro que tenemos una noción instintiva y racional de la gravedad, entonces,
en mi opinión, está igual de claro que la gravedad también ejerce una atracción
irresistible sobre nuestra imaginación. Precisamente porque estamos inmersos en un
campo gravitatorio y sujetos a él en cada instante de nuestra vida, nos fascinan de
manera natural esos escenarios en los que la gravedad es débil o incluso está ausente.
¿Quién no ha deseado alguna vez lanzarse desde un acantilado o desde la cima de una
montaña y… volar? ¿Quién no ha imaginado que es un astronauta a bordo de la Estación
Espacial Internacional o un personaje de una película de ciencia ficción que flota sin
esfuerzo de un lugar a otro en gravedad cero? A mí me pasa a menudo… Dicho de otro
modo, la gravedad atrae nuestra atención. Estimula nuestra imaginación precisamente
porque es la única «fuerza» de la cual somos conscientes y porque sabemos muy bien lo
difícil que es escapar de ella. ¿Acaso no fue la imaginación lo que llevó, primero a
Isaac Newton\index{Newton!Isaac} y posteriormente a Albert Einstein\index{Einstein!Albert},
a explicar, de maneras muy distintas, las leyes que gobiernan la gravedad?

% DOCX paragraph 87
Hay muchos ejemplos que ilustran el poderoso tirón de la gravedad sobre nuestra
imaginación, pero me voy a limitar a uno, que me parece representativo y fácil de
entender. En 2013, el director español Alfonso
Cuarón\index{Cuaron, Alfonso (""Gravity"")@Cuarón, Alfonso (""Gravity"")}
rodó una película con el emblemático título \emph{Gravity}. Durante casi dos horas,
la película no habla de otra cosa que de la gravedad o, más bien, de su ausencia.
Sin embargo, poca gente sabe que \emph{Gravity} batió el récord de taquilla del primer
fin de semana entre las películas estrenadas aquel otoño. Hay quienes sostienen que
este éxito se debe exclusivamente a los actores principales, las dos superestrellas
de Hollywood Sandra Bullock y George Clooney. En mi opinión, sin embargo, desempeñó
un papel fundamental el hecho de que, nos guste o no, no podemos escapar de la gravedad
ni de su irresistible tirón sobre nuestra mente.\index{conocimiento!imaginativo}

% DOCX paragraph 90
\section[Una de cuatro, pero diferente]{Una de cuatro, pero decididamente diferente de las otras}\label{sec:01-interactions}

% DOCX paragraph 92
Las consideraciones hechas hasta ahora proporcionan un punto de partida excelente
para presentar otra importante reflexión sobre la función que tiene la gravedad en
nuestra comprensión de la naturaleza. La física moderna nos enseña que hay cuatro
tipos de interacciones fundamentales, que esencialmente describen todos los procesos
que tienen lugar en el universo: la \emph{interacción electromagnética}\index{interaccion@interacción!electromagnetica@electromagnética},
la \emph{interacción fuerte}\index{interaccion@interacción!fuerte}, la
\emph{interacción débil}\index{interaccion@interacción!debil@débil} y la
\emph{interacción gravitatoria}\index{interaccion@interacción!gravitatoria}.

% DOCX paragraph 93
La primera, la interacción electromagnética, entre otras cosas, permite leer este libro
con independencia del formato que se esté empleando. De hecho, las ondas
electromagnéticas\index{ondas!electromagneticas@electromagnéticas}
---fotones\index{foton@fotón} o, sencillamente, luz--- emanan de la página que el lector
tiene ante sí y alcanzan sus ojos. Entonces, esas ondas se convierten en señales
eléctricas transmitidas por el nervio óptico al cerebro, el cual, gracias a una compleja
combinación de intercambios eléctricos y químicos, las traduce a las palabras que aquí
se leen. La interacción electromagnética es también responsable de la cohesión y la
dinámica de las moléculas de las que estamos compuestos. Sin ella, ni siquiera
existiríamos como seres humanos; nuestras moléculas se desperdigarían como papelitos
movidos por el viento. La teoría que describe la interacción electromagnética es bien
conocida, tanto en la física clásica, donde se expresa mediante las ecuaciones de
Maxwell, como en la física cuántica cuando nos ocupamos de las partículas elementales,
contexto en el que también se la conoce como teoría de la \emph{electrodinámica cuántica}.

% DOCX paragraph 94
En cambio, la segunda, la interacción fuerte, se desarrolla en la escala más pequeña
de la naturaleza que nos es accesible: del orden de unos pocos fermi, también llamados
femtómetros, es decir, unas pocas milésimas de una milmillonésima de milímetro. La
interacción fuerte es cien veces más intensa que la interacción electromagnética y
se ejerce entre los \emph{quarks}, los constituyentes de partículas elementales como
los protones o los neutrones. En realidad, esta interacción también está presente a
una escala un poco mayor, es decir, dentro de los núcleos
atómicos\index{atomico@atómico!nucleo@núcleo}, que generalmente tienen dimensiones del
orden de 10 fermi, donde adopta el nombre de \emph{fuerza nuclear fuerte}. En el caso
de los protones y neutrones, las partículas mediadoras de la interacción fuerte se
llaman \emph{gluones}\index{gluones}, mientras que en la fuerza nuclear fuerte están
representadas por los \emph{piones}\index{piones}. Básicamente, la interacción fuerte
actúa como un «pegamento» que mantiene unidos los núcleos atómicos, desde los átomos
más pequeños, como el hidrógeno, hasta los más grandes, como el uranio. También regula
lo que ocurre cuando dos protones chocan entre sí a una velocidad casi igual a la de
la luz o cuando nace una estrella de neutrones\index{estrella de neutrones}, fenómeno
que estudiaremos detalladamente en el capítulo 5. La teoría que describe esta
interacción está bien desarrollada y se llama \emph{cromodinámica
cuántica}\index{cromodinamica cuantica@cromodinámica cuántica}. Lamentablemente, dada
la complejidad de la teoría y de las ecuaciones que la describen, a menudo es difícil
hacer predicciones precisas, sobre todo si las energías son elevadas o si el número
de partículas que intervienen es grande, como en el caso de las estrellas de neutrones.

% DOCX paragraph 95
La tercera interacción, la débil, es responsable de la desintegración radiactiva de
algunos núcleos atómicos y actúa entre \emph{leptones}\index{leptones} ---una clase
de partículas de la que forman parte los electrones\index{electron@electrón}, sin
duda los más conocidos del grupo--- y quarks. Gracias a esta interacción, los
neutrinos\index{neutrino}, partículas muy ligeras producidas por materia de alta
densidad y temperatura, como la del centro del Sol\index{Sol}, interactúan muy poco
(«débilmente») con protones y neutrones, la materia ordinaria de la que estamos
compuestos. Conviene recordar que, en este preciso momento, miles de millones de
neutrinos atraviesan nuestro cuerpo: neutrinos emitidos por el Sol hace unos ocho
minutos y que han viajado hasta nosotros a una velocidad próxima a la de la luz.
Si no sentimos que estas partículas nos atraviesan, es precisamente porque solo
interactúan débilmente con la materia ordinaria o \emph{hadrónica}\index{materia hadronica (ordinaria)@materia hadrónica (ordinaria)}
de la que estamos hechos. En otras palabras, no hay por qué preocuparse: somos
esencialmente transparentes para los neutrinos. La teoría que describe la interacción
débil también está muy bien desarrollada y se puede combinar con la del
electromagnetismo, de modo que, en realidad, hablamos de la \emph{interacción electrodébil}.

% DOCX paragraph 96
Y así, por último, llegamos a la cuarta interacción, la gravitatoria. Aunque todavía
no ha llegado el momento de explicar con detalle en qué consiste y cómo se relaciona
con uno de los conceptos más sutiles y elegantes de la física teórica ---la curvatura
del espaciotiempo---, podemos reflexionar ya sobre lo que la distingue de las demás.
La gravedad es la única interacción física de la que somos conscientes. Es la única
cuya influencia sobre nuestro cuerpo percibimos ahora mismo, mientras leemos este
libro. De hecho, tanto si estamos acostados en una cama, hundidos en un sillón o de
pie, \emph{sabemos} que alguna «fuerza» afecta a nuestra posición. Sin ella, estaríamos
levitando, flotando libremente como un astronauta en la Estación Espacial Internacional.

% DOCX paragraph 97
Este punto es esencial para alcanzar una comprensión cabal del título del capítulo
y captar su significado más profundo: la gravedad atrae nuestra atención porque
podemos experimentar su existencia de forma directa y tangible, a diferencia de lo
que ocurre con las otras interacciones fundamentales. En cambio, no podemos apreciar
el nivel de cohesión de las moléculas que conforman nuestro cuerpo, ni la escasa
frecuencia con la que interactuamos con un neutrino, ni el número de partículas
radiactivas que emitimos.\textsuperscript{\hyperlink{note:01-01}{1}}

% DOCX paragraph 98
El hecho de que la podamos experimentar directa y conscientemente basta para otorgar
a la gravedad un lugar especial entre las interacciones fundamentales y, en mi
opinión, situarla por encima de todas las demás. Lo que la hace aún más peculiar
es que esta conciencia nos acompaña cada segundo de nuestra vida, incluso antes
de nacer. Incluso de forma inconsciente, sabemos que la gravedad existe mucho antes
de rasguñarnos las rodillas cuando aprendemos a caminar y correr, y desde luego
antes de encontrarnos con las leyes de la física en la escuela o en la universidad.

% DOCX paragraph 101
\section{Pero ¿qué es la gravedad y cómo funciona?}\label{sec:01-gravity}

% DOCX paragraph 103
Es muy probable que muchos de los lectores de este libro estén convencidos de poder
dar una respuesta razonable a estas sencillas preguntas: «¿Qué es la gravedad y cómo
funciona?». Esto se debe a que la conocemos tanto instintiva como racionalmente y a
que, en la escuela o en la universidad, nos dieron una explicación «culta». Ahora
bien, es igual de probable que las explicaciones que recibimos no fueran del todo
correctas, aunque tampoco estuvieran completamente traídas por los pelos. Quiero
decir que lo que nos dijeron o enseñaron no es realmente \emph{incorrecto}, pero
tampoco es \emph{correcto}.

% DOCX paragraph 104
La razón de esta aparente contradicción es que la gravedad se puede entender en
distintos niveles. Por ejemplo, la representación de la gravedad propuesta por
Newton\index{Newton!Isaac} es más sencilla y más intuitiva: implica la existencia
de una \emph{fuerza gravitatoria}\index{fuerza!gravitatoria} y su
descripción matemática es relativamente simple. Al mismo tiempo, es posible una
comprensión diferente y más profunda: la que sugirió Einstein, que implica una
visión geométrica del espacio y el tiempo y una descripción matemática mucho
más compleja.

% DOCX paragraph 105
En los capítulos que siguen nos espera una especie de evolución: desarrollaremos
nuestra comprensión de la gravedad. La primera fase de esta evolución ---y, por
tanto, el primer destino de nuestro viaje virtual--- nos conducirá desde la
comprensión elemental de la gravedad tal como está codificada en nuestro cerebro,
incluso a nivel instintivo, hasta la descripción que ofrece la teoría de
Newton\index{Newton!Isaac}. En cambio, la segunda fase ---o el siguiente destino
del viaje--- nos llevará a la descripción, matemáticamente compacta y físicamente
profunda, que propuso Einstein con su teoría de la relatividad
general\index{relatividad!general|see{relatividad general}}\index{relatividad general},
plasmada con elegancia en sus ecuaciones de campo.

% DOCX paragraph 106
Así pues, aprenderemos a conectar lo que sabemos y experimentamos en nuestro planeta,
donde los campos gravitatorios son muy débiles, con lo que observamos en el universo,
donde campos gravitatorios muchísimo más fuertes dan vida a fenómenos maravillosos
como los agujeros negros, las estrellas de neutrones y las ondas
gravitatorias\index{ondas!gravitatorias (vease ondas gravitatorias)@gravitatorias (véase ondas gravitatorias)}.
De esta forma nos daremos cuenta de que en nuestro conocimiento y nuestra comprensión
de la gravedad influye en gran medida la manera en que esta se manifiesta en nuestro
planeta. Ahora bien, también nos daremos cuenta de que tenemos que abandonar esa
visión por ser incorrecta y porque restringe demasiado nuestra imaginación.

% DOCX paragraph 107
En conjunto, quedará clara la respuesta correcta a las preguntas planteadas:
«La gravedad es simplemente la manifestación de la curvatura del espaciotiempo».
En este momento se trata todavía de una afirmación críptica, lo sé. Pero, tal como
se menciona en la introducción, hace falta una cierta dosis de paciencia para
embarcarse en el viaje que nos espera. Puedo asegurar que todo quedará mucho más
claro al final del capítulo 3.
